\subsection{A topology forced on every surviving countable simple target}

\begin{claim}{A.29}{Proved}{Unimodular closure and a direct countable topology construction}
Let $F$ be a nonabelian simple group. If finitely many
words $w_i(x)\in F*\langle x\rangle$ satisfy
\[
 w_i(1)\ne1\quad\text{for all }i,
 \qquad
 F-\{1\}=\bigcup_i\{g\in F:w_i(g)=1\},
 \tag{A.29.1}
\]
then an explicitly constructed one-variable equation of
total exponent $1$ has no solution in $F$.

Consequently every countable nonabelian simple group that
solves every one-variable equation of total exponent
$1$ or $-1$ admits a nondiscrete Hausdorff metrizable
group topology. More generally, the same topological
conclusion holds for any countably infinite group $F$
with the following finite-avoidance property:
\[
 \begin{gathered}
 p_1(1),\ldots,p_r(1)\ne1,
 \quad p_i\in F*\langle x\rangle
 \\
 \Longrightarrow\quad
 \text{there is }a\ne1\text{ with }p_i(a)\ne1
 \text{ for every }i.
 \end{gathered}
 \tag{A.29.2}
\]
In particular every surviving simple countable quotient
of the $S^2$ example in A.28 has such a topology.
\end{claim}
\begin{proof}
We use $[a,b]=aba^{-1}b^{-1}$. Suppose first that
(A.29.1) holds, and write the list as $w_1,\ldots,w_m$.
Put $v_1=w_1$. Given $v_{i-1}(1)\ne1$, there is
$c_i\in F$ such that
\[
 [v_{i-1}(1),c_iw_i(1)c_i^{-1}]\ne1.
 \tag{A.29.3}
\]
Otherwise $v_{i-1}(1)$ centralizes every conjugate of
$w_i(1)$ and thus its normal closure, which is $F$.
This contradicts the triviality of the center of a
nonabelian simple group. Define
\[
 v_i(x)=[v_{i-1}(x),c_iw_i(x)c_i^{-1}].
 \tag{A.29.4}
\]
Then $v_i(1)\ne1$. If any of
$w_1(g),\ldots,w_i(g)$ is $1$, the recursive formula
gives $v_i(g)=1$. Thus $v_m(1)\ne1$ and
$v_m(g)=1$ for every $g\ne1$.

When $m\geq2$ the word $v_m$ has exponent zero,
since it is a commutator. When $m=1$, choose $c$
with $[w_1(1),c]\ne1$ and replace $v_1$ by
$[w_1(x),c]$. The resulting word $v$ in either case
has exponent zero, is nonidentity at $1$, and
vanishes at every other element. Therefore
\[
 u(x)=v(x)x
 \tag{A.29.5}
\]
has exponent $1$, but $u(1)=v(1)\ne1$ and
$u(g)=g\ne1$ for $g\ne1$. The equation $u(x)=1$
has no solution.

It follows that a nonabelian simple group solving
all the specified unimodular equations satisfies
(A.29.2). Such a group is infinite: for a finite
one the words $xg^{-1}$, indexed by $g\ne1$,
would give exactly the finite cover (A.29.1).
It remains to prove the topological assertion
directly from (A.29.2).

\paragraph{An auxiliary proper length.}
Enumerate the nonidentity elements of $F$ as
$f_1,f_2,\ldots$. Give each symbol $f_i$ and
$f_i^{-1}$ cost $i$, and let $\ell(g)$ be the
minimum cost of a finite word representing $g$,
with $\ell(1)=0$. This minimum exists because
the costs are nonnegative integers and every
element has a representing symbol. It satisfies
\[
 \ell(g^{-1})=\ell(g),\qquad
 \ell(gh)\leq\ell(g)+\ell(h),\qquad
 \ell(g)\geq1\quad(g\ne1).
 \tag{A.29.6}
\]
Each finite-radius ball for $\ell$ is finite.
Indeed a word of cost at most $R$ uses only
the finitely many symbols of cost at most $R$
and has at most $R$ letters. Put
\[
 c(g)=2^{\ell(g)},\qquad c(kg)\leq c(k)c(g).
 \tag{A.29.7}
\]

\paragraph{The finite recursive constraints.}
We choose nonidentity elements $a_1,a_2,\ldots$.
A formal letter of level $i$ is a triple
$(g,i,\varepsilon)$, where $g\in F$ and
$\varepsilon\in\{1,-1\}$. Its value and cost are
\[
 ga_i^{\varepsilon}g^{-1},\qquad
 c(g)2^{-i},
 \tag{A.29.8}
\]
respectively. At stage $n$ impose this invariant:
for every $j\leq n$, no finite product of formal
letters of levels at most $n$ and total cost
less than $2^{-j}$ has value $f_j$.
The empty product, of value $1$, is harmless.
For $j=n$ the condition is automatic, since
each nonempty product costs at least $2^{-n}$.

Suppose the earlier choices satisfy the invariant
at stage $n-1$. For each $j<n$ list all products
of formal letters of levels at most $n$ and cost
less than $2^{-j}$. There are finitely many.
Every occurring letter obeys
\[
 c(g)2^{-i}<2^{-j},\qquad
 \ell(g)<i-j\leq n-j,
 \tag{A.29.9}
\]
so the possible conjugators belong to a finite
union of finite $\ell$-balls. Every letter
costs at least $2^{-n}$, so each listed product
has fewer than $2^{n-j}$ letters.

In each listed product replace $a_n$ by the
free variable $x$ and append $f_j^{-1}$.
This gives finitely many words $p(x)$ over $F$.
At $x=1$ the level-$n$ letters disappear.
The remaining earlier product still has cost
less than $2^{-j}$ and, by the induction
invariant, is not $f_j$. Thus every such
$p$ satisfies $p(1)\ne1$. Property (A.29.2)
supplies $a_n\ne1$ at which all these words
are nonidentity. This proves the invariant
at stage $n$. At stage $1$ any $a_1\ne1$
works, since there are no nonautomatic constraints.

\paragraph{The resulting metric.}
Take a symmetric collection of weighted
letters consisting of all the formal letters
(A.29.8), together with every nonidentity
$g\in F$ as an ordinary letter of cost
$\ell(g)$. Distinct formal letters can have
the same value; they retain their own costs.
Let $L(h)$ be the infimum of the total costs
of finite products of these letters evaluating
to $h$. Ordinary letters make $L$ finite
everywhere. Inverting a product preserves its
cost, and concatenating products adds costs,
so
\[
 L(1)=0,\qquad L(h^{-1})=L(h),\qquad
 L(hk)\leq L(h)+L(k).
 \tag{A.29.10}
\]
For each $j$ one moreover has
\[
 L(f_j)\geq2^{-j}>0.
 \tag{A.29.11}
\]
To see this, a product costing less than
$2^{-j}$ cannot contain an ordinary nonidentity
letter, whose cost is at least $1$. If its
maximum formal level is $n<j$, every letter
already costs more than $2^{-j}$, so only the
empty product is possible. If $n\geq j$, the
stage-$n$ invariant excludes value $f_j$.
This proves (A.29.11).

Consequently
\[
 d(g,h)=L(g^{-1}h)
 \tag{A.29.12}
\]
is a genuine left-invariant metric. The letter
of level $i$ with conjugator $1$ shows
$L(a_i)\leq2^{-i}$. Thus the nonidentity
$a_i$ converge to $1$, and this metric
topology is nondiscrete.

\paragraph{Continuity of the group operations.}
For fixed $k\in F$ put
\[
 D_k=\max\{c(k),1+2\ell(k)\}.
 \tag{A.29.13}
\]
Conjugating a formal letter with conjugator
$g$ by $k$ produces the formal letter with
conjugator $kg$. Its cost is at most $c(k)$
times the original cost by (A.29.7).
Conjugating an ordinary nonidentity letter
$g$ gives an ordinary nonidentity letter of
cost at most
\[
 \ell(kgk^{-1})\leq2\ell(k)+\ell(g)
       \leq(1+2\ell(k))\ell(g).
 \tag{A.29.14}
\]
Conjugate any representing product and take
infima. This proves
\[
 L(khk^{-1})\leq D_k L(h)\qquad(h\in F).
 \tag{A.29.15}
\]
Since
$(gh)^{-1}(g'h')=h^{-1}(g^{-1}g')h\,h^{-1}h'$,
we obtain
\[
 \begin{split}
 d(gh,g'h')&\leq D_{h^{-1}}d(g,g')+d(h,h'),\\
 d(g^{-1},g'^{-1})&\leq D_g d(g,g').
 \end{split}
 \tag{A.29.16}
\]
These inequalities prove continuity of
multiplication at each pair and of inversion
at each point. The metric therefore defines
the claimed Hausdorff group topology.

The simple quotients remaining in A.28
solve every one-variable unimodular equation
by its explicit quotient-lifting argument.
The result just proved applies to each such
countable quotient.
\end{proof}

This additional topological requirement does
not decide the remaining quotient class.
The construction provides no assertion of
completeness or compactness, and no equation
system without a solution has been inferred
from the topology alone.
